% Generated draft: requires mathematical and source review.
\documentclass[11pt]{article}
\usepackage[margin=25mm]{geometry}
\usepackage{fontspec}
\setmainfont{Latin Modern Roman}
\usepackage{amsmath,amssymb,mathtools,unicode-math}
\setmathfont{Latin Modern Math}
\usepackage{enumitem,needspace,xurl,hyperref}
\setlength{\parindent}{0pt}
\setlength{\parskip}{5pt}
\setlength{\emergencystretch}{3em}
\newcommand{\N}{\mathbb{N}}
\newcommand{\Z}{\mathbb{Z}}
\newcommand{\Q}{\mathbb{Q}}
\newcommand{\R}{\mathbb{R}}
\newcommand{\C}{\mathbb{C}}
\newcommand{\F}{\mathbb{F}}
\newcommand{\A}{\mathbb{A}}
\newcommand{\PP}{\mathbb{P}}
\newcommand{\E}{\mathbb{E}}
\newcommand{\eps}{\varepsilon}
\newcommand{\dd}{\,\mathrm{d}}
\begin{document}
\section*{UnsolvedMath ID: 30005610}
\textbf{Source problem code:} catalogue entry
\subsection{Definitions and mathematical statement}
The supplied record poses three related questions. It does not specify a complete class of domains, admissible Robin coefficients, or conventions for eigenfunctions at multiple eigenvalues. The following definitions give a conditional formulation in a standard setting; they do not resolve these omissions or assert that this setting is the full intended scope.

For this conditional formulation, let \(\Omega\subset\R^d\), \(d\geq2\), be a bounded connected domain with smooth boundary, and let \(\nu\) be its outward unit normal. For a boundary condition \(B\), consider real Laplace eigenfunctions satisfying
\[
-\Delta u_k^B=\lambda_k^B u_k^B\quad\text{in }\Omega.
\]
List eigenvalues in nondecreasing order, with multiplicity, and fix a real orthonormal eigenbasis \(\mathcal U_B=(u_k^B)_{k\geq1}\). The two boundary conditions relevant here are
\[
B=D:\quad u|_{\partial\Omega}=0,
\qquad
B=R_\beta:\quad \partial_\nu u+\beta u=0
\quad\text{on }\partial\Omega,
\]
where \(\beta\in\R\) is fixed in this reference formulation. Whether the source intends constant parameters, variable boundary coefficients, or another admissible class is unspecified.

For a nonzero real eigenfunction \(u\), let \(\mu_\Omega(u)\) be the number of connected components of
\(\{x\in\Omega:u(x)\ne0\}\). Define the actual nodal-count limit superior by
\[
P(\Omega,B;\mathcal U_B)
:=\limsup_{k\to\infty}\frac{\mu_\Omega(u_k^B)}{k}.
\]
This expresses the source's intended Pleijel constant as an actual limit superior, rather than merely an upper bound. The displayed fraction interprets the source's corrupted expression ``lim sup Nkk'' as \(\limsup_{k\to\infty}N_k/k\). The eigenbasis is displayed explicitly because the source does not specify how eigenfunctions in multiple eigenspaces are selected or whether any optimization over them is intended.

\begin{enumerate}
\item Determine whether the Robin and Dirichlet Pleijel constants on the same domain always agree. In the conditional setting above, the proposed equality is
\[
P(\Omega,R_\beta;\mathcal U_{R_\beta})
\stackrel{?}{=}
P(\Omega,D;\mathcal U_D).
\]
The word ``always'' requires the intended domain class, admissible Robin data, and quantifiers over eigenfunction choices to be supplied before this becomes a uniquely specified universal assertion.

\item Define the harmonic Dirichlet-to-Neumann operator by
\[
\Lambda f=\left.\partial_\nu Hf\right|_{\partial\Omega},
\qquad
\Delta(Hf)=0\text{ in }\Omega,
\qquad
\left.Hf\right|_{\partial\Omega}=f,
\]
initially for smooth boundary data, with its usual self-adjoint realization in \(L^2(\partial\Omega)\). A real Steklov eigenfunction \(v_j\) and its boundary trace \(f_j\) satisfy
\[
\Delta v_j=0\text{ in }\Omega,
\qquad
\partial_\nu v_j=\sigma_jv_j\text{ on }\partial\Omega,
\qquad
f_j=v_j|_{\partial\Omega},
\qquad
\Lambda f_j=\sigma_j f_j.
\]
Order \((\sigma_j)_{j\geq1}\) nondecreasingly with multiplicity, including zero. For a nonzero real boundary eigenfunction, define
\[
\mu_{\partial\Omega}(f_j)
:=\#\{\text{connected components of }
\{x\in\partial\Omega:f_j(x)\ne0\}\}.
\]
The question asks for information or bounds on these boundary nodal counts. It specifies no particular bound, asymptotic rate, or dependence on \(j\), \(\sigma_j\), or the geometry. The source motivates this question by observing that \(\Lambda^2\) has the same principal symbol as the local boundary Laplacian: in the smooth reference setting, both principal symbols are \(|\xi|_{g_{\partial\Omega}}^2\). This observation does not identify the two operators or specify a nodal-count conclusion.

\item Find further consequences of Steklov--Robin duality. With the sign convention above, the basic correspondence concerns the same nonzero function satisfying
\[
\Delta v=0\text{ in }\Omega,
\qquad
\partial_\nu v=\sigma v\text{ on }\partial\Omega,
\]
and the Robin problem with boundary parameter \(\beta=-\sigma\) at interior spectral value zero. The source requests additional consequences but gives no particular conclusion to prove and does not specify whether a more general spectral-parameter version of this duality is intended.
\end{enumerate}

The source labels these questions open. Its supplied dated assessment reports generic equality for balls while describing the general-domain equality, general Dirichlet-to-Neumann nodal-count questions, and further duality consequences as unresolved. Neither the meaning of ``generically'' nor the hypotheses of the reported ball result are supplied, so that assertion is retained as qualified source context rather than formulated as a theorem.
\subsection{Short English statement}
The questions ask whether actual Robin and Dirichlet nodal-count limit superiors agree on a fixed domain, what bounds hold for nodal domains of Steklov boundary traces, and what further consequences follow from Steklov--Robin duality. The source leaves domain hypotheses, Robin data, multiplicity conventions, and the desired bounds or consequences unspecified. The displayed formulation is conditional and requires human mathematical review.
\subsection{Sources}
\begin{itemize}
\item Local UnsolvedMath registry: supplied complete source record, permanent ID \(30005610\), title \emph{Pleijel Constants and Steklov--Robin Duality}, exact source problem code \texttt{catalogue entry}, in the Oberwolfach Reports Open Problems collection. No individual proposer is named; the proposed year is recorded as 2023.
\item Supplied report citation: \emph{Geometric Spectral Theory (2024)}, Oberwolfach Reports, pp.~2017--2106, DOI \texttt{10.4171/owr/2023/36}. \url{https://doi.org/10.4171/owr/2023/36}. The record describes the report as documenting workshop discussions and retains the quoted passage as the context for the extracted questions. The citation's 2024 designation and its 2023 DOI component are preserved without reconciliation.
\item Supplied literature source, labeled ``Hassannezhad primary article'': \url{https://doi.org/10.1112/blms.13005}. The record supplies no article title, theorem number, or detailed hypotheses for the claimed generic result for balls.
\item Supplied literature source, labeled ``Oberwolfach Report 36'': \url{https://doi.org/10.4171/owr/2023/36}.
\item Literature assessment dated exactly \texttt{2026-08-22}: the supplied curation labels the problem open and reports that Robin Pleijel-constant equality is proved generically for balls but remains open for general domains, with general Dirichlet-to-Neumann nodal counts and further Steklov--Robin consequences likewise unresolved. This is supplied dated curation derived from the OWR record, not an independently verified literature assessment; its time-sensitive claims require independent checking.
\item The registry marks the statement \texttt{corrected\_verified} and states that it was reconstructed and checked against the official report. These are supplied provenance claims, not verification performed for this draft. Authoritative references are supplied, but their contents have not been independently inspected here.
\end{itemize}
\end{document}
